\section{Árboles Trinomiales con Mallas Adaptativas}
\label{sec:impl_arboles_mallas}

\subsection{Descripción del Algoritmo}
\label{subsec:algoritmo_arboles_mallas}

La implementación del Modelo de Malla Adaptativa (AMM) para opciones de barrera difiere fundamentalmente de los árboles trinomiales estándar en su fase de inicialización. Mientras que en un árbol estándar se define arbitrariamente el número de pasos de tiempo ($N$) y se derivan los parámetros espaciales, en el AMM la geometría de la malla está dictada por las condiciones de contorno (la barrera) y la profundidad de refinamiento deseada.
El algoritmo comienza con una fase de calibración estricta para asegurar que los nodos de la malla fina coincidan exactamente con la barrera del precio, evitando así el ``problema del límite''.

\subsubsection{Calibración de la Geometría (Malla Gruesa)}

La construcción de la malla base (denotada como conjunto de nodos $A$ o malla gruesa) requiere resolver un sistema de restricciones acopladas para determinar el paso de precio ($h$) y el paso de tiempo ($k$).

\paragraph{Restricción Espacial: Alineación con la Barrera} 

Para valorar con precisión una opción barrera, es imperativo que la barrera física $H$ coincida con una capa horizontal de nodos en el árbol. En el AMM, definimos la resolución en función de la distancia entre el precio inicial del activo ($S_0$) y la barrera ($H$).

Sea $M$ el número de niveles de refinamiento de la malla fina (donde $M=0$ implica un árbol trinomial estándar restringido). El algoritmo debe calcular el paso de precio de la malla gruesa, $h$, de tal manera que, al dividirlo por $2M$, el paso más fino cubra exactamente la distancia a la barrera.

La fórmula de inicialización para $h$ es: 
\begin{equation} 
  h = 2^M \times | \ln(S_0) - \ln(H) | \label{eq:amm_h_calib} 
\end{equation}

Esta ecuación garantiza que el nodo $B_{1,0}$ (en la malla fina) coincida con el precio inicial $S_0$ y que el nodo inferior de la malla fina coincida con $H$.

\textbf{Nota de Implementación:} El árbol se construye en el espacio logarítmico $X=ln(S)$. Por tanto, la barrera se convierte en una constante $ln(H)$ para todos los pasos de tiempo, simplificando la verificación de knock-out.

\paragraph{Restricción Temporal: El Parámetro de Estiramiento} 

Una vez fijado $h$ por la restricción espacial, el paso de tiempo $k$ no puede ser elegido libremente. Debe cumplir dos condiciones simultáneas:
\begin{enumerate} 
  \item Mantener la estabilidad numérica y la convergencia de la varianza (típicamente gobernado por un parámetro libre $\lambda\approx3$).
  
  \item Resultar en un número entero de pasos de tiempo $N$ hasta el vencimiento $T$.
\end{enumerate}

Dado que calcular $k$ directamente a partir de $h$ y la volatilidad $\sigma$ raramente produce un entero para $T$, se implementa el ``Parámetro de Estiramiento'' (\textit{Stretch Parameter}). El algoritmo calcula un $k$ provisional y luego lo "estira" ligeramente para ajustar $N$ al entero más cercano.
La fórmula implementada, derivada de la Ecuación 17 de Figlewski y Gao (1999), es:
\begin{equation} 
  k = \frac{T}{\text{int}\left( \frac{j \sigma^2 T}{h^2} \right)} \label{eq:amm_k_stretch} 
\end{equation}

Donde: 
\begin{itemize} 
  \item $j$: Parámetro de estabilidad (generalmente se establece $j=3$ para minimizar el error de distribución). 
  
  \item $int(\cdot)$: Función que toma la parte entera. 
\end{itemize}

El denominador representa el número total de pasos de tiempo ($N$) de la malla gruesa. Este ajuste asegura que el árbol termine exactamente en $T$ sin necesidad de truncar o interpolar el último paso.

\paragraph{Cálculo de Probabilidades de la Malla Base} 

Con $h$ y $k$ calibrados, se calculan las probabilidades de transición para los nodos de la malla gruesa. A diferencia de los árboles centrados en la media, acá el árbol está anclado en la barrera, por lo que incluimos el término \textit{drift} en las probabilidades.
Definimos el \textit{drift} $\alpha$ como: 
\begin{equation} 
  \alpha = r - q - \frac{\sigma^2}{2} 
\end{equation} 

Donde $r$ es la tasa libre de riesgo y $q$ la tasa de dividendos.

Las probabilidades de movimiento hacia arriba ($p_u$), al medio ($p_m$) y hacia abajo ($p_d$) para la malla gruesa se implementan como:

\begin{equation} 
  p_u = \frac{1}{2} \left( \frac{\sigma^2 k}{h^2} + \frac{\alpha^2 k^2}{h^2} + \frac{\alpha k}{h} \right) 
\end{equation}
\begin{equation} 
  p_d = \frac{1}{2} \left( \frac{\sigma^2 k}{h^2} + \frac{\alpha^2 k^2}{h^2} - \frac{\alpha k}{h} \right) 
\end{equation}
\begin{equation} 
  p_m = 1 - p_u - p_d 
\end{equation}

\textbf{Validación:} El código debe incluir una aserción (\textit{assert}) inmediatamente después de este cálculo para verificar que $p_u,p_m,p_d\ge0$. Si el parámetro de estiramiento $j$ es cercano a $3$, las probabilidades serán positivas y estables.

\subsubsection{Gestión de la Interfaz entre Mallas}

Una vez calibrada y calculada la malla gruesa (Nodos del conjunto $A$), el siguiente paso es construir la malla fina (Nodos del conjunto $B$) en la región crítica cercana a la barrera.

Para una opción barrera, el flujo de información es descendente (Malla Gruesa → Malla Fina). Esto significa que utilizaremos los valores ya calculados en los nodos $A$ para interpolar los valores de los nodos frontera de $B$.

El desafío de implementación radica en los ``Nodos Intermedios'' (o nodos colgantes). Dado que la malla fina tiene una resolución temporal $4$ veces mayor que la gruesa, existen nodos $B$ en los tiempos $t+k/4$, $t+2k/4$ y $t+3k/4$ que no tienen un nodo $A$ correspondiente en el mismo instante temporal. Para valorarlos, debemos proyectarlos hacia los nodos $A$ del siguiente paso de tiempo completo ($t+k$).

\paragraph{Cálculo de Probabilidades Ajustadas (Ecuación de Transición)} 

Consideremos los nodos $B_{2,j}$ (ubicados un paso de precio grueso por encima de la barrera). Estos nodos sirven como el ``techo'' de la malla fina y deben conectarse a los nodos de la malla gruesa $A$.

Para conectar un nodo intermedio ubicado en el instante $t+\tau$ con los nodos $A$ en el instante $t+k$, debemos calcular probabilidades trinomiales usando un paso de tiempo efectivo ($\Delta t_{eff}$) que cubra la distancia restante hasta la columna de nodos $A$.

Las probabilidades ajustadas para un nodo con paso de tiempo efectivo $\Delta t_{eff}$ se implementan según las siguientes fórmulas (derivadas de la Ec. 12 de Figlewski y Gao):

\begin{equation} 
  p_u(\Delta t_{eff}) = \frac{1}{2} \left( \frac{\sigma^2 \Delta t_{eff}}{h^2} + \frac{\alpha^2 (\Delta t_{eff})^2}{h^2} + \frac{\alpha \Delta t_{eff}}{h} \right) 
\end{equation}

\begin{equation} 
  p_d(\Delta t_{eff}) = \frac{1}{2} \left( \frac{\sigma^2 \Delta t_{eff}}{h^2} + \frac{\alpha^2 (\Delta t_{eff})^2}{h^2} - \frac{\alpha \Delta t_{eff}}{h} \right) 
\end{equation}
\begin{equation} 
  p_m(\Delta t_{eff}) = 1 - p_u - p_d 
\end{equation}

\textbf{Lógica de Bucle:} En el código, para cada intervalo grueso entre $t$ y $t+k$, se deben calcular tres sets de probabilidades distintos para los nodos frontera: 

\begin{itemize} 
  \item \textbf{Para nodos en $t+k/4$:} Usar $\Delta t_{eff}=3k/4$. 

  \item \textbf{Para nodos en t+2k/4:} Usar $\Delta t_{eff}=2k/4$. 
  
  \item \textbf{Para nodos en t+3k/4:} Usar $\Delta t_{eff}=k/4$. 
\end{itemize}

\paragraph{Valoración de los Nodos Frontera} 

Con las probabilidades ajustadas, calculamos el valor de la opción en estos nodos intermedios descontando desde los nodos $A$ del futuro.

Sea $V(B_{2,j})$ el valor de un nodo en la frontera superior de la malla fina. Su valor se obtiene de los nodos de la malla gruesa $A_u$ (arriba), $A_m$ (medio) y $A_d$ (abajo) correspondientes al tiempo $t+k$:

\begin{equation} 
  V(B_{2,j}) = e^{-r\Delta t_{eff}} \left[ p_u(\Delta t_{eff}) V(A_u) + p_m(\Delta t_{eff}) V(A_m) + p_d(\Delta t_{eff}) V(A_d) \right] \label{eq:amm_interface_val} 
\end{equation}

\textbf{Nota crítica de implementación:} Observe que el factor de descuento es dinámico y debe ser $e^{-r\Delta t_{eff}}$ Dado que las probabilidades de transición se calculan proyectando el nodo intermedio directamente hacia los nodos de la malla gruesa ($A$), el descuento debe aplicarse sobre toda la distancia temporal efectiva ($\Delta t_{eff}$) que los separa, y no por un único micro-paso de $k/4$.

\paragraph{Relleno de la Malla Fina} 

Una vez calculados los valores en el ``techo'' de la malla fina ($B_{2,j}$) y sabiendo que el ``piso'' de la malla es la barrera ($B_{0,j}=0$ o valor de Rebate), procedemos a calcular los nodos internos ($B_{1,j}$) que se encuentran a $h/2$ de la barrera.

Estos nodos no requieren conexión con la malla gruesa ya que se valoran mediante inducción hacia atrás estándar dentro de la propia malla fina.

\begin{enumerate} 
  \item \textbf{Geometría Fina:} Usar paso de precio $h_{fine}=h/2$ y paso de tiempo $k_{fine}=k/4$.
  
  \item \textbf{Probabilidades Finas:} Calcular $p_u^\prime,p_m^\prime, p_d^\prime$ usando las fórmulas estándar trinomiales con $h_{fine}$ y $k_{fine}$.

  \item \textbf{Recursión:} 
  \begin{equation} 
    V(B_{1,j}) = e^{-r(k/4)} \left[ p'u V(B{2, j+1}) + p'm V(B{1, j+1}) + p'd V(B{0, j+1}) \right] 
  \end{equation}
\end{enumerate}

Este proceso se repite retrocediendo desde el vencimiento $T$ (o desde el final del bloque injertado) hasta llegar a $t=0$, obteniendo finalmente el valor en el nodo $B_{1,0}$, que corresponde al precio exacto de la opción con alta precisión.

\subsubsection{Valoración Backward y Recursividad}

Una vez establecidas las probabilidades de transición y las reglas de interfaz, el proceso de valoración se ejecuta mediante un algoritmo secuencial estricto. A diferencia de un árbol estándar que se resuelve en un solo barrido desde $T$ hasta $0$, el AMM requiere una valoración por etapas jerárquicas.

\paragraph{Fase 1: Resolución de la Malla Gruesa (Conjunto $A$)} 

El algoritmo comienza resolviendo completamente la malla base, ignorando temporalmente la existencia de la malla fina. 

\begin{enumerate} 
  \item Se inicializan los valores terminales en $t=T$ para todos los nodos $A_{i,N}$ basándose en el payoff de la opción: $max(S-K,0)$ o $max(K-S,0)$. 
  
  \item Se ejecuta la inducción hacia atrás (\textit{backward induction}) estándar hasta $t=0$, aplicando la condición de barrera en cada paso: si un nodo $A_{i,j}$ cae por debajo de la barrera (en el caso de down-and-out), su valor se fuerza a $0$ (o al valor del Rebate). 
  
  \item Resultado: Al finalizar esta fase, tenemos una matriz de valores $V(A_{i,j})$ que servirá como fuente de datos (``oráculo'') para los niveles inferiores. 
\end{enumerate}

\paragraph{Fase 2: Inyección de Valores en la Frontera Fina} 

Antes de resolver la malla fina, debemos poblar sus bordes utilizando los datos de la Fase 1. El código debe iterar sobre cada paso de tiempo grueso $j$ y calcular los valores de los nodos de interfaz $B_{2,*}$ (el ``techo'' de la malla fina).

El bucle debe manejar la interpolación temporal descrita en la sección anterior: 
\begin{itemize} 
  \item Para los nodos $B$ que coinciden con tiempos $A$ ($t,t+k,\dots$), se copia el valor directamente desde la malla gruesa. La calibración de la Ec. \ref{eq:amm_h_calib} garantiza que los nodos espaciales estén perfectamente alineados, por lo que no se requiere ningún tipo de interpolación espacial. 
  
  \item Para los nodos intermedios ($t+k/4,t+2k/4,\dots$), se ejecuta la función de proyección usando la Ecuación \ref{eq:amm_interface_val}. 
\end{itemize}

\paragraph{Fase 3: Resolución de la Malla Fina (Conjunto $B$)} 

Con las condiciones de frontera superior ($B_{2,j}$ calculados en Fase 2) e inferior ($B_{0,j}=0$ por la barrera), se resuelve el interior de la malla fina. 

\begin{enumerate} 
  \item Se comienza en el vencimiento $T$ (o en el punto donde la malla fina se une a la gruesa). 
  
  \item Se retrocede paso a paso con $\Delta t=k/4$. 
  
  \item En cada paso, se calcula el valor de los nodos internos $B_{1,j}$ utilizando las probabilidades trinomiales estándar de la malla fina y los valores de los tres nodos conectados en el paso $j+1$. 
\end{enumerate}

\paragraph{Generalización Recursiva (Para $M>1$)} 

Si el modelo utiliza múltiples niveles de refinamiento ($M>1$), este proceso se vuelve recursivo. La propiedad de isomorfismo del AMM permite que, una vez resuelta la malla de Nivel 1 (Nodos $B$), esta actúe como la ``Malla Gruesa'' para el Nivel 2 (Nodos $C$).
Un código preliminar es:

\begin{verbatim}
def Resolver_AMM(S0, K, T, r, q, sigma, H, M):
  # 1. CALIBRACIÓN GLOBAL (Para la malla más profunda)
  h_coarse_base, k_base, N_base = calibrar_geometria(S0, H, T, sigma, M)

  # 2. CONSTRUIR MALLA BASE (NIVEL 0)
  A_grid = build_coarse_mesh(S0, K, T, r, q, sigma, H, h_coarse_base, k_base, N_base)
  
  # 3. BUCLE DE REFINAMIENTO (NIVEL 1 hasta M)
  current_grid = A_grid
  current_h, current_k, current_N = h_coarse_base, k_base, N_base
  
  for level in range(1, M + 1):
    fine_grid = build_fine_mesh(current_grid, current_h, current_k, current_N, r, 
                  q, sigma, H, K) # La nueva malla fina se "injerta" en la actual
    
    current_grid = fine_grid # La fina de hoy es la gruesa de mañana
    current_h, current_k, current_N = actualizar_parametros(...)
  
  return current_grid[1, 0] # Retornar el valor actual (t=0) del nodo central

def build_coarse_mesh(S0, K, T, r, q, sigma, H, h, k, N):
  pu_A, pm_A, pd_A = calcular_probs(k, h, r, q, sigma) # Eq. 9 para malla gruesa
  
  A_grid = initialize_coarse_mesh()
  A_grid[max_price_nodes, N] = actualizar_payoff_terminal(...)
  A_grid[:, 0:N-1] = Standard_Backward_Induction(...)
  return A_grid

def build_fine_mesh(coarse_grid, h_coarse, k_coarse, N_coarse, r, q, sigma, H, K):   
  total_fine_steps = N_coarse * 4
  B_grid = initialize_fine_mesh()
  
  # 1. INYECCIÓN DE NODOS (Eq. 11)
  for j in range(N_coarse): # Lógica de interpolación (Eq. 12 y 13)
    t_start = j * 4
    B_grid[2, t_start] = coarse_grid[1, j] # El source siempre es la malla previa
    V_Ad, V_Am, V_Au = coarse_grid[:, j+1] # Nodos intermedios (Interpolación)
    
    for sub_step in range(1, 4):
      dt_eff = (4 - sub_step) * (k_coarse / 4) # Distancia efectiva a próximo A
      pu_adj, pm_adj, pd_adj = calcular_probs(dt_eff, h_coarse, r, q, sigma)
      
      # Valorar usando los nodos A del FUTURO
      B_grid[2, t_start + sub_step] = (pu_adj * V_Au + 
                                        pm_adj * V_Am + 
                                        pd_adj * V_Ad) * np.exp(-r * dt_eff)
    
  # 2. RELLENO INTERNO (Backward Induction)
  h_fine, k_fine = h_coarse / 2, k_coarse / 4
  pu, pm, pd = calcular_probs(k_fine, h_fine, r, q, sigma)

  B_grid[2, total_fine_steps] = coarse_grid[1, N_coarse]
  B_grid[:, total_fine_steps] = actualizar_payoff_terminal(...)
  B_grid[:, 0:total_fine_steps-1] = Standard_Backward_Induction(...)
  return B_grid
\end{verbatim}

Esta estructura garantiza que la información fluya correctamente desde la resolución más baja hacia la más alta (flujo downward para barreras), permitiendo obtener el precio de alta precisión en $t=0$.

\subsection{Detalles de Implementación}
\label{sec:detalles_impl_arboles}

\subsubsection{Estructuras de Datos y Memoria}
\label{sec:data_structures}

Uno de los errores más comunes al implementar métodos de refinamiento de malla es instanciar una matriz global de alta resolución y dejar la mayoría de sus entradas vacías o inutilizadas. La eficiencia del AMM radica en que la malla fina es un ``parche'' local y no una estructura global. Por tanto, la gestión de memoria debe tratar las mallas como entidades separadas pero vinculadas.

\paragraph{Arquitectura de Almacenamiento Disperso} 

En lugar de crear una matriz de tamaño ($4N\times 4N$) que consumiría memoria cuadrática, la implementación debe utilizar dos estructuras distintas:

\begin{enumerate} 
  \item Matriz Base ($A$): Un array bidimensional de tamaño $(N+1)\times(2N+1)$ para almacenar los valores de la malla gruesa.
  
  \item Matriz de Banda ($B$): Una estructura optimizada que almacena solo los nodos necesarios a lo largo de la barrera.
\end{enumerate}

Para un modelo de barrera con un nivel de refinamiento ($M=1$), la malla fina no cubre todo el espacio de precios, sino solo una franja estrecha adyacente a $H$. Según el análisis de complejidad, la cantidad de nodos adicionales es lineal respecto al tiempo: 

\begin{equation} 
  Nodos_{adicionales} = 10 \times N 
\end{equation} 

Donde $N$ es el número de pasos de tiempo de la malla gruesa.

\textbf{Recomendación de Implementación:} Representar la malla fina $B$ como una matriz de dimensiones ($3\times 4N$). 

\begin{itemize} 
  \item La dimensión vertical es $3$ porque solo necesitamos calcular los nodos en los niveles $h/2$ (interno), $h$ (frontera superior) y $0$ (barrera/frontera inferior). 
  
  \item La dimensión horizontal es $4N$ porque la resolución temporal es cuatro veces mayor. 
\end{itemize} 
Esto reduce el consumo de memoria drásticamente en comparación con refinar todo el árbol.

\paragraph{Complejidad Computacional y Escalabilidad} 

La ventaja crítica del AMM debe documentarse en esta sección. Al añadir niveles de profundidad, el número de nodos crece, pero de manera controlada.

Para un nivel de profundidad $m$, el número de nodos nuevos añadidos por cada paso de tiempo grueso es: 

\begin{equation} 
  \Delta Nodos(m) = 10 \times 4^{m-1} 
\end{equation}

Esto implica que para un árbol estándar de $N=100$ pasos (aprox. $10,000$ nodos totales): 

\begin{itemize} 
  \item \textbf{AMM Nivel 1:} Añade solo 1,000 nodos (incremento del $10\%$). $Total \approx 11,200 nodos$. 
  
  \item \textbf{Equivalente Estándar (Refinado Global):} Para obtener la misma resolución espacial cerca de la barrera sin AMM, se requeriría cuadruplicar $N$ a $400$. Esto generaría aprox. $160,000 nodos$. 
\end{itemize}

El código debe reflejar esta eficiencia: el tiempo de ejecución para el AMM debe ser comparable al del árbol trinomial base, mientras que el error de convergencia se reduce cuadráticamente.

\subsubsection{Código Relevante}
\label{sec:codigo_relevante}

La implementación del AMM requiere traducir las ecuaciones de probabilidad ajustada y la lógica recursiva en estructuras de control eficientes. A continuación se presentan los tres componentes críticos que permiten la convergencia estable del modelo.
\paragraph{1. Calibración Geométrica y Parámetro de Estiramiento} 

Para garantizar la estabilidad numérica descrita por Figlewski y Gao (1999), el paso de tiempo $k$ no puede ser arbitrario; debe derivarse del paso de precio $h$ para mantener el parámetro de estiramiento $j\approx3$. El siguiente snippet muestra cómo se fuerza esta restricción para obtener un número entero de pasos $N$.

\begin{lstlisting}[language=Python, caption={Cálculo de pasos de tiempo enteros (Eq. 17)}] 
  def calibrar_geometria(S0, H, T, sigma, M): 
    distancia = |log(S0) - log(H)|

    # Eq. 18: El paso de precio crece exponencialmente con M
    h_coarse = (2**M) * distancia

    # Eq. 11 y 17: Ajuste del paso de tiempo para j=3
    # Se redondea N al entero mas cercano para evitar pasos fraccionarios
    N_ideal = (T * 3 * sigma**2) / (h_coarse**2)
    N_steps = int(N_ideal)

    k_coarse = T / N_steps
    return h_coarse, k_coarse, N_steps
\end{lstlisting}

\paragraph{2. Probabilidades de Transición en la Interfaz} 

La mayor complejidad del AMM reside en los nodos ``colgantes'' de la interfaz (Fila 2 de la malla fina). A diferencia de los nodos estándar, estos requieren probabilidades que varían según su distancia temporal al siguiente nodo grueso, tal como se define en la Ecuación (12).

\begin{lstlisting}[language=Python, caption={Cálculo dinámico de probabilidades para nodos de interfaz}] 
  def calcular_probs(dt_eff, dx, r, q, sigma): 
    """ 
    Calcula probabilidades para sub-pasos k/4, 2k/4, 3k/4. dt_eff: 
      Tiempo restante hasta el proximo nodo grueso (T, T-k/4...) 
    """ 
    drift = r - q - 0.5 * sigma**2

    # Terminos de la Ecuacion (12)
    term1 = (sigma**2 * dt_eff) / (dx**2)
    term2 = (drift**2 * dt_eff**2) / (dx**2)
    term3 = (drift * dt_eff) / dx

    pu = 0.5 * (term1 + term2 + term3)
    pd = 0.5 * (term1 + term2 - term3)
    pm = 1.0 - pu - pd

    return pu, pm, pd
\end{lstlisting}

\paragraph{3. El Bucle de Refinamiento Recursivo} 

\subsection{Validación} 
\label{sec:validacion_arboles}

Para verificar la corrección de la implementación y demostrar la superioridad del AMM sobre los árboles estándar, se realizaron pruebas de convergencia comparando los resultados del modelo contra la solución analítica cerrada para opciones de barrera continua.

\subsubsection{Metodología de Validación (Benchmark)} 

Es crucial notar que el AMM converge a la solución de tiempo continuo, no a la discreta. Por lo tanto, el benchmark correcto no es un árbol binomial de pasos altos, sino la fórmula de Merton (1973) para opciones down-and-out call:

\begin{equation} 
  C_{teo} = C_{BS}(S_0, K) - \left( \frac{H}{S_0} \right)^{2\lambda} C_{BS}(H^2/S_0, K) \label{eq:merton_benchmark} 
\end{equation} 

Donde $\lambda=\frac {r-q-\sigma^2/2} {\sigma^2}$ y $C_{BS}$ es la fórmula estándar de Black-Scholes.

\subsubsection{Prueba de Estrés: El Problema del Límite} 

Se evaluó el desempeño del modelo en el escenario crítico donde el precio inicial $S_0$ está extremadamente cerca de la barrera, una situación donde los métodos tradicionales fallan por falta de resolución o inestabilidad numérica.

\textbf{Configuración del Experimento:} 

\begin{itemize} 
  \item Barrera $H=90$, Strike $K=100$, $T=1$, $\sigma=0.25$, $r=0.1$. 
  
  \item Precio Inicial Crítico: $S_0=90.125$ (Distancia de $0.14\%$ a la barrera). 
\end{itemize}

\textbf{Resultados Comparativos:}

\begin{table}[h] \centering 
  \begin{tabular}{|l|c|c|c|} 
    \hline \textbf{Método} & \textbf{Pasos de Tiempo (N)} & \textbf{Valor Calculado} & \textbf{Error Relativo} \\ 
    \hline Solución Exacta (Merton) & - & \textbf{$\approx 0.161648$} & - \\ 
    \hline Árbol Trinomial ($M=0$) & N/A & \textit{Fallo (Oscilación)} & $>100\%$ \\ 
    \hline AMM Nivel 2 ($M=2$) & $6083$ & $0.161312$ & $0.21\%$ \\ 
    \hline AMM Nivel 4 ($M=4$) & $380$ & \textbf{$0.161532$} & $0.07\%$ \\ 
    \hline 
  \end{tabular} \caption{Comparación de convergencia cerca de la barrera ($S_0=90.125$).} \label{tab:amm_results} 
\end{table}

\textbf{Análisis:} Como se observa en la Tabla \ref{tab:amm_results}, el árbol estándar es incapaz de valorar la opción con una malla razonable porque el paso de precio mínimo es mayor que la distancia $S_0-H$. El AMM con $M=4$ logra una precisión de $4$ decimales manteniendo el costo computacional bajo ($N=100$ pasos base), resolviendo efectivamente el problema del límite mediante la inserción local de nodos.

\subsubsection{Análisis de Eficiencia} 

Las pruebas de tiempo de ejecución confirman que el costo computacional del AMM escala linealmente. Añadir un nivel de refinamiento incrementa el tiempo de cálculo en aproximadamente un factor constante pequeño (correspondiente a valorar la franja de 3 filas), en contraste con el factor de $\times16$ ($4^2$) que requeriría refinar todo el árbol globalmente. Esto valida la eficiencia del método para aplicaciones en tiempo real.
